\section{Proof of the main theorem}\label{sec:conclusion}

We assemble the preceding constructions, preserving the specified
isomorphisms of volume lines.

Start with the smooth compact K\"ahler model \(S\) prepared in
Section~\ref{sec:prelim}. Corollary~\ref{orders:bdescent} proves that
\[
 M_{S_1}=\nu^*M_S
\]
for every smooth compact K\"ahler modification \(\nu:S_1\to S\).
This proves part~\textup{(a)} of Theorem~\ref{thm:main}.

For part~\textup{(b)}, Proposition~\ref{spread:product} supplies a
proper surjection \(r:P\to S\) from a smooth compact connected
auxiliary base, and an actual product decomposition of log-plurivolume
lines on a dense open \(P^\circ\). The connected deepest-stratum
selection in that proposition is compatible with the original
top root line by Proposition~\ref{hodge:residue}.
Theorem~\ref{orders:trace} therefore identifies the extension of
the original factor on the left with \(r^*(mM_S)\), after taking
the necessary common multiples and testing on local unipotent
covers.

By Proposition~\ref{compare:semiample}, the projective-factor line
has a semiample extension, with the required logarithmic norm
bounds. By Proposition~\ref{periods:semiample}, so do the remaining
highest-form factor lines after a further finite cover and
resolution of \(P\). Replace \(P\) by that compact smooth model.
It is still proper and surjective over \(S\).
All powers and all finitely many finite covers can be chosen once
for this family. Their composite and a common further exponent
suffice simultaneously for every factor.

The product isomorphism on \(P^\circ\) is induced by the actual
volume maps, and its norms agree up to the fixed covering-degree
constant. Corollary~\ref{orders:residuenorm} and the two factor
propositions supply upper and lower powers-of-logarithm estimates
for local extended frames on general transverse disks.
Lemma~\ref{descent:unit}, applied also to the inverse isomorphism,
extends the product identity to all of \(P\).
Thus, for some positive integer \(a\), the actual line
\(r^*\cO_S(aM_S)\) is a tensor product of semiample lines.
A common power is globally generated.

Lemma~\ref{descent:semiample} now gives a positive multiple of
\(aM_S\) generated at every point of \(S\).
The exponent is independent of the point, and the conclusion
holds for the holomorphic line bundle itself.
This proves part~\textup{(b)} and completes
Theorem~\ref{thm:main}.
\qed

\begin{remark}
The compact auxiliary base is used only to establish generation.
It is not asserted to be the modification in
Theorem~\ref{thm:main}. The latter is the compact K\"ahler model
\(S\) on which the threshold calculation already proves pullback
compatibility. Nor is a choice of global divisor representatives
for canonical or moduli line bundles on the original nonprojective
manifolds needed anywhere in the argument.
\end{remark}
